\chapter*{存疑八条}
\phantomsection
\addcontentsline{toc}{chapter}{存疑八条}
\markboth{存疑八条}{存疑八条}

\noindent
下列八片，或文不成句，或无法归类，或明显为后世插入。本书照录其字，不作解释。

下编不处理这八条。理由见导论第五节。

\begin{cunyi}{存疑一　（原作 12 号）}
□□□□，其四□□。凡如是者，不问。
\end{cunyi}
\xiaojì{可辨者九字。虽有「不问」二字，无以知其不问者为何。}

\begin{cunyi}{存疑二　（原作 29 号）}
角。声。□。此三者，非一。
\end{cunyi}
\xiaojì{「角」「声」二字清楚，第三个字蚀尽。三字之间关系不明。}

\begin{cunyi}{存疑三　（原作 36 号）}
凡我不与。凡我不与。
\end{cunyi}
\xiaojì{同句重刻两遍，字口深浅略异，或为两人所刻。}

\begin{cunyi}{存疑四　（原作 44 号）}
三十九。四十一。四十三。
\end{cunyi}
\xiaojì{三个奇数连续，无单位，无动词。可解为编号，亦可解为计量。}

\begin{cunyi}{存疑五　（原作 55 号）}
……以其故，故不可。其故不可言。
\end{cunyi}
\xiaojì{首尾俱失，中段完整。文气与其余各片不类，疑为转录他人之语。}

\begin{cunyi}{存疑六　（原作 62 号）}
一人。二人。三人。皆不至。
\end{cunyi}
\xiaojì{「皆不至」，所至何处，不明。}

\begin{cunyi}{存疑七　（原作 68 号）}
此句与上句同。此句与上句同。
\end{cunyi}
\xiaojì{亦为重刻，与存疑三重出而不同形，二者关系不明。}

\begin{cunyi}{存疑八　（原作 88 号）}
□□□。其下三字，吾不能辨。辨者当自为之。
\end{cunyi}
\xiaojì{「辨者当自为之」，径对后来者说话，与全书体例不合。此片真伪存疑。}
